\documentclass[11pt,a4paper]{article}
\usepackage[margin=2.5cm]{geometry}
\usepackage{amsmath,amssymb}
\usepackage{enumitem}
\usepackage{xcolor}
\usepackage{hyperref}
\usepackage{tikz}

\definecolor{popblue}{HTML}{1F77B4}
\definecolor{sampred}{HTML}{D90012}
\definecolor{paramgreen}{HTML}{2CA02C}

\hypersetup{colorlinks=true, linkcolor=popblue, urlcolor=popblue}

\title{\textbf{The Connection Between Pearson Correlation\\and the Regression Slope}}
\author{Metric Academy}
\date{}

\begin{document}
\maketitle

\section{Setup: Two Familiar Quantities}

Suppose we have $n$ data points $(x_1, y_1), \dots, (x_n, y_n)$.

\subsection*{The Pearson Correlation Coefficient}

The Pearson correlation $r$ measures the \emph{strength and direction} of the
linear relationship between $x$ and $y$:

\[
r = \frac{\sum_{i=1}^{n}(x_i - \bar{x})(y_i - \bar{y})}
         {\sqrt{\sum_{i=1}^{n}(x_i - \bar{x})^2} \;\cdot\;
          \sqrt{\sum_{i=1}^{n}(y_i - \bar{y})^2}}
\]

Key properties:
\begin{itemize}[nosep]
  \item Always between $-1$ and $+1$.
  \item Unitless --- it does not depend on the scale of $x$ or $y$.
  \item $r = +1$: perfect positive linear relationship.
  \item $r = -1$: perfect negative linear relationship.
  \item $r = 0$: no \emph{linear} relationship (there could still be a nonlinear one).
\end{itemize}

\subsection*{The Regression Slope (from the Normal Equation)}

In simple linear regression we fit $\hat{y} = \hat{\beta}_0 + \hat{\beta}_1 x$.
The normal equation gives us the slope:

\[
\hat{\beta}_1
= \frac{\sum_{i=1}^{n}(x_i - \bar{x})(y_i - \bar{y})}
       {\sum_{i=1}^{n}(x_i - \bar{x})^2}
\]

The slope tells us: \emph{for each one-unit increase in $x$, how much does
$\hat{y}$ change?} Unlike $r$, the slope carries units (e.g.\ kg per cm).

%-------------------------------------------------------------------
\section{The Key Relationship}

Look at the two formulas above. They share the same numerator
$\sum(x_i - \bar{x})(y_i - \bar{y})$. The difference is only in the
denominator. With a bit of algebra we can express one in terms of the other.

\begin{center}
\fcolorbox{popblue}{blue!5}{%
  \parbox{0.75\textwidth}{\centering\vspace{6pt}
    \textbf{The slope is the correlation, rescaled by the ratio of standard
    deviations:}
    \[
      \boxed{\;\hat{\beta}_1 \;=\; r \;\cdot\; \frac{s_y}{s_x}\;}
    \]
    where $\;s_x = \sqrt{\frac{1}{n-1}\sum(x_i-\bar{x})^2}\;$ and
    $\;s_y = \sqrt{\frac{1}{n-1}\sum(y_i-\bar{y})^2}$.
  \vspace{6pt}}%
}
\end{center}

\subsection*{Derivation}

Start from the slope formula and multiply top and bottom by
$\frac{1}{n-1}$:

\begin{align*}
\hat{\beta}_1
&= \frac{\sum(x_i-\bar{x})(y_i-\bar{y})}{\sum(x_i-\bar{x})^2}
 = \frac{\tfrac{1}{n-1}\sum(x_i-\bar{x})(y_i-\bar{y})}
        {\tfrac{1}{n-1}\sum(x_i-\bar{x})^2}
 = \frac{s_{xy}}{s_x^2}
\end{align*}

where $s_{xy} = \frac{1}{n-1}\sum(x_i-\bar{x})(y_i-\bar{y})$ is the sample
covariance. Now recall that $r = \frac{s_{xy}}{s_x \, s_y}$, so
$s_{xy} = r \, s_x \, s_y$. Substituting:

\[
\hat{\beta}_1
= \frac{r \, s_x \, s_y}{s_x^2}
= r \cdot \frac{s_y}{s_x}
\]

That's it --- one line of algebra connects them. \hfill $\square$

%-------------------------------------------------------------------
\section{What Does This Tell Us?}

\subsection{1. Same sign, always}

Because $s_x > 0$ and $s_y > 0$, the sign of $\hat{\beta}_1$ equals the sign
of $r$. A positive correlation always means an upward-sloping regression line,
and vice versa. This is a sanity check you can use in practice.

\subsection{2. Correlation is the ``standardized slope''}

Imagine you standardize both variables before running regression:
\[
z_x = \frac{x - \bar{x}}{s_x}, \qquad
z_y = \frac{y - \bar{y}}{s_y}
\]

Now both variables have mean 0 and standard deviation 1, so
$s_{z_y}/s_{z_x} = 1$. Therefore:
\[
\hat{\beta}_1^{\text{(standardized)}}
= r \cdot \frac{s_{z_y}}{s_{z_x}}
= r \cdot \frac{1}{1}
= r
\]

\textbf{The Pearson correlation is exactly the slope you get when both
variables are on the same scale.} This is why $r$ is unitless and bounded
between $-1$ and $+1$: it strips away the units and scale differences,
leaving only the strength and direction of the linear association.

\subsection{3. The $R^2$ connection}

In simple (one-predictor) regression:
\[
R^2 = r^2
\]

So squaring the Pearson correlation gives you the \emph{coefficient of
determination} --- the fraction of variance in $y$ explained by the linear
relationship with $x$.

For example, if $r = 0.8$, then $R^2 = 0.64$, meaning 64\% of the
variability in $y$ is accounted for by $x$.

\subsection{4. Units are the only real difference}

\begin{center}
\begin{tabular}{lcc}
\hline
& \textbf{Slope} $\hat{\beta}_1$ & \textbf{Correlation} $r$ \\
\hline
Units        & units of $y$ per unit of $x$ & unitless \\
Range        & $(-\infty, +\infty)$ & $[-1, +1]$ \\
Measures     & rate of change & strength of association \\
Comparable across datasets? & No & Yes \\
\hline
\end{tabular}
\end{center}

The factor $s_y / s_x$ is what converts between the two. It rescales the
unitless correlation into a slope with real-world units.

%-------------------------------------------------------------------
\section{A Concrete Example}

Suppose we are studying the relationship between hours studied ($x$) and exam
score ($y$) for 50 students.

\begin{itemize}[nosep]
  \item $\bar{x} = 5$ hours, \quad $s_x = 2$ hours
  \item $\bar{y} = 72$ points, \quad $s_y = 10$ points
  \item $r = 0.85$
\end{itemize}

Then the regression slope is:
\[
\hat{\beta}_1 = 0.85 \times \frac{10}{2} = 4.25 \;\text{points per hour}
\]

\textbf{Interpretation:}
\begin{itemize}[nosep]
  \item The correlation $r = 0.85$ tells us there is a strong positive linear
        association.
  \item The slope $\hat{\beta}_1 = 4.25$ tells us that each additional hour of
        studying is associated with about 4.25 extra points on the exam.
  \item $R^2 = 0.85^2 = 0.72$, so 72\% of the variance in exam scores is
        explained by hours studied.
\end{itemize}

Both numbers describe the same linear relationship, just in different
``languages''.

%-------------------------------------------------------------------
\section{Summary}

\begin{enumerate}
  \item The regression slope and Pearson correlation encode the \emph{same}
        linear relationship.
  \item They differ only by a scale factor: \;
        $\hat{\beta}_1 = r \cdot \dfrac{s_y}{s_x}$.
  \item Correlation is the standardized version of the slope --- what you get
        when both variables are on the same scale.
  \item They always share the same sign.
  \item In simple regression, $R^2 = r^2$.
\end{enumerate}

\vfill
\noindent\rule{\textwidth}{0.4pt}\\[4pt]
{\small
\textbf{Sources:}
\begin{itemize}[nosep,leftmargin=*]
  \item Penn State STAT 501 --- Pearson Correlation Coefficient:
        \url{https://online.stat.psu.edu/stat501/lesson/1/1.6}
  \item Sebastian Raschka --- Pearson R vs.\ Simple Linear Regression:
        \url{https://sebastianraschka.com/faq/docs/pearson-r-vs-linear-regr.html}
  \item UT Austin --- Pearson Correlation and Linear Regression:
        \url{https://sites.utexas.edu/sos/guided/inferential/numeric/bivariate/cor/}
\end{itemize}
}

\end{document}
